Infinity is not a real number to which ordinary cancellation rules apply. The symbols $\infty-\infty$ and $\dfrac{\infty}{\infty}$ do not specify a numerical limit. For example, as $t\to\infty$, both terms of each difference tend to infinity, yet $t-t=0$, $(t+1)-t=1$, and $2t-t=t\to\infty$. Similarly $\dfrac{t}{t}=1$, $\dfrac{2t}{t}=2$, and $\dfrac{t^2}{t}=t\to\infty$. A limit depends on the full expressions. Moreover $\dfrac{1}{t}\to0$ means convergence, not that a positive finite reciprocal equals zero. The proposed argument uses undefined operations and cannot establish $1=0$.
